\begin{table}[htbp]\centering\footnotesize
\caption{Las cuentas de la declaratoria urbana contra el archivo de las 43.636 AGEB}\label{cua:zapcuentas}
\begin{tabular}{>{\raggedright\arraybackslash}p{5.0cm}>{\raggedleft\arraybackslash}p{2.2cm}>{\raggedleft\arraybackslash}p{2.0cm}>{\raggedleft\arraybackslash}p{2.2cm}}
\toprule
Cuenta & Declarado en el anexo K & Sobre la clave actual & Sobre las claves del archivo \\
\midrule
AGEB urbanas prioritarias & 43.636 & 43.636 & 43.636 \\
Entidades federativas & 32 & 32 & 32 \\
\textbf{Municipios} & 2.423 & 2.423 & \textbf{2.416} \\
\textbf{Localidades urbanas} & 4.531 & 4.531 & \textbf{4.540} \\
\midrule AGEB dentro de los 1.575 municipios ZAP rurales & no se publica & 25.836 & --- \\
\bottomrule
\end{tabular}
\\[2pt]\begin{minipage}{0.92\textwidth}\scriptsize Fuente: anexo K de la Gaceta 7121 para lo declarado; y \texttt{2027/2027\_zap-urbanas-agebs\_variables.zip}, descargado el 10 de septiembre de 2026, para lo calculado. \textbf{Los dos renglones en negritas son el aviso}: municipios y localidades sólo cierran contra la última columna del archivo, \emph{clave de localidad actual a junio de 2026}. Quien cuente sobre las columnas de clave que el archivo trae al frente publicará 2.416 y 4.540. El traslape con la lista rural sí lo publica la nota metodológica: el archivo lo reproduce al AGEB.\end{minipage}
\end{table}
